\chapter{Linear and Circular Polarization, Operators, and Expectation Values}

\section{The Smallest Useful Quantum System}

Take a photon moving in the \(z\)-direction with definite momentum and set
aside its translational motion. What remains is its transverse polarization.
Classically we picture an electric field oscillating along an axis in the
\(xy\)-plane. For a single photon that transverse degree of freedom becomes
a particularly clean quantum system: the state space is only
two-dimensional, but it already contains superposition, incompatible
measurements, complex amplitudes, and probabilistic outcomes.

Begin with an ideal linear polarizer. The conducting direction of a
wire-grid polarizer supports a current and therefore suppresses the electric
field component parallel to the wires. The perpendicular component is
transmitted. Consequently the \emph{transmission axis}, rather than the wire
direction, labels the polarizer.

The same apparatus has two roles. If a photon emerges from a polarizer, the
apparatus has prepared a state. If a prepared photon is sent into a second
polarizer, that apparatus tests the state. Equal preparation and test axes
give certain transmission; orthogonal axes give certain rejection.

\begin{classroomqa}
\audiencequestion{Could ordinary classical particles imitate this first
experiment?\lecturetimestamp{00:09:01--00:09:25}}
\lecturerresponse{Yes. One could imagine bullets carrying little flags and
a mechanical slit that passes one flag orientation while blocking the
other. Nothing nonclassical has appeared while the only preparations and
tests are the same two orthogonal directions.}
\end{classroomqa}

The difficulty begins with rotational symmetry. A polarizer can be turned
to any angle, so a photon can be prepared along a continuum of linear
polarization axes. A fixed analyzer nevertheless has only two outcomes:
the photon passes or it does not. No single photon divides into a
transmitted fraction and a rejected fraction. Continuous preparations and
discrete outcomes must coexist.

This is the central puzzle. A classical wave divides continuously, and the
flagged-bullet model remains deterministic. A photon at an intermediate
angle instead gives one indivisible event whose probability depends on the
relative orientation of the apparatus.

Mathematically this polarization problem is isomorphic to a spin-one-half
system: both use a two-dimensional complex state space and noncommuting
two-outcome observables. Polarization axes are somewhat easier to picture,
so they provide the cleaner geometry for the present discussion.

\section{The Axial Basis and Its Observable}

Choose horizontal and vertical transmission axes and call them \(x\) and
\(y\). Linear polarization is axial: horizontal-to-the-right is not
different from horizontal-to-the-left. Reversing the drawn arrow leaves the
same polarization axis.

\begin{figure}[tbp]
  \centering
  \includegraphics[width=0.38\textwidth]{lecture_06_figure_02.png}
  \caption{The double-headed mark records a polarization axis without
  assigning it a preferred direction.}
  \lectureframe{00:14:32}
\end{figure}

Represent the two orthogonal states by
\begin{equation}
  \ket{x}=
  \begin{pmatrix}1\\0\end{pmatrix},
  \qquad
  \ket{y}=
  \begin{pmatrix}0\\1\end{pmatrix}.
  \label{eq:l6-xy-basis}
\end{equation}
They obey
\begin{equation}
  \braket{x}{x}=\braket{y}{y}=1,
  \qquad
  \braket{x}{y}=0.
  \label{eq:l6-xy-orthogonality}
\end{equation}
These are two basis vectors, not the only two state vectors. Every
normalized superposition
\begin{equation}
  \ket{\psi}=a\ket{x}+b\ket{y},
  \qquad
  |a|^2+|b|^2=1,
  \label{eq:l6-general-two-state}
\end{equation}
is a possible pure polarization state. Multiplication of the entire vector
by one common phase does not change the physical state.

\begin{classroomqa}
\audiencequestion{Are the entries \(1,0\) and \(0,1\) completely
arbitrary?\lecturetimestamp{00:16:50--00:18:38}}
\lecturerresponse{The coordinate convention is chosen with the axes, but
the vectors must be normalized and mutually orthogonal. Those conditions
encode total probability one and perfect distinguishability. Their common
overall phase is irrelevant, so several column vectors can represent the
same physical ray.}
\end{classroomqa}

To turn a test into an observable, attach numbers to its outcomes. Record
\(+1\) when an \(x\)-analyzer transmits the photon and \(-1\) when it rejects
it. The labels are conventional; \(7\) and \(12\) would distinguish the
outcomes too. The symmetric choice \(\pm1\) gives
\begin{equation}
  \hat P_0\ket{x}=+\ket{x},
  \qquad
  \hat P_0\ket{y}=-\ket{y},
  \label{eq:l6-P0-eigenvalues}
\end{equation}
and therefore
\begin{equation}
  \hat P_0=
  \begin{pmatrix}
    1&0\\
    0&-1
  \end{pmatrix}
  =\ket{x}\bra{x}-\ket{y}\bra{y}.
  \label{eq:l6-P0}
\end{equation}

\begin{figure}[tbp]
  \centering
  \includegraphics[width=0.58\textwidth]{lecture_06_figure_03.png}
  \caption{Orthogonality and the first matrix action in the
  \((\ket{x},\ket{y})\) basis.}
  \lectureframe{00:22:00}
\end{figure}

The operator can be constructed from known eigenstates and recorded
values, as we have just done. In other problems the order is reversed:
one begins with an operator and solves for its eigenvectors and eigenvalues.
For example, the wave-number operator
\(-i\,d/dx\) is given first, and its plane-wave eigenstates are then found.
Either route expresses the same postulate: an observable has definite value
\(\lambda\) in a state satisfying
\(\hat A\ket{\lambda}=\lambda\ket{\lambda}\).

\section{A Second Basis at Forty-Five Degrees}

Turn the preparation polarizer through \(45^\circ\). The transmitted state
has equal \(x\)- and \(y\)-components:
\begin{equation}
  \ket{d_+}
  =\frac{\ket{x}+\ket{y}}{\sqrt2}
  =\frac{1}{\sqrt2}
    \begin{pmatrix}1\\1\end{pmatrix}.
  \label{eq:l6-diagonal-plus}
\end{equation}
Its orthogonal companion, the backslash polarization, can be chosen as
\begin{equation}
  \ket{d_-}
  =\frac{\ket{x}-\ket{y}}{\sqrt2}
  =\frac{1}{\sqrt2}
    \begin{pmatrix}1\\-1\end{pmatrix}.
  \label{eq:l6-diagonal-minus}
\end{equation}
Putting the minus sign on the other component multiplies the whole state by
\(-1\) and changes no physics. Directly,
\begin{equation}
  \braket{d_-}{d_+}=\frac12(1-1)=0.
  \label{eq:l6-diagonal-orthogonality}
\end{equation}

The amplitude that a slash-polarized photon passes an \(x\)-analyzer is
\begin{equation}
  \braket{x}{d_+}=\frac{1}{\sqrt2}.
\end{equation}
The probability is the absolute square,
\begin{equation}
  \operatorname{Prob}(x|d_+)
  =|\braket{x}{d_+}|^2=\frac12,
  \qquad
  \operatorname{Prob}(y|d_+)=\frac12.
  \label{eq:l6-diagonal-probabilities}
\end{equation}
The amplitude and the probability are not interchangeable. Amplitudes may
be negative or complex; probabilities are nonnegative real numbers.

Now attach \(+1\) to the slash outcome and \(-1\) to the backslash outcome.
The corresponding observable is
\begin{equation}
  \hat P_{\pi/4}
  =\ket{d_+}\bra{d_+}-\ket{d_-}\bra{d_-}
  =
  \begin{pmatrix}
    0&1\\
    1&0
  \end{pmatrix}.
  \label{eq:l6-P45}
\end{equation}
Exchanging which diagonal outcome is called \(+1\) simply replaces this
operator by \(-\hat P_{\pi/4}\).
Indeed,
\begin{equation}
  \hat P_{\pi/4}
  \begin{pmatrix}1\\1\end{pmatrix}
  =
  \begin{pmatrix}1\\1\end{pmatrix},
  \qquad
  \hat P_{\pi/4}
  \begin{pmatrix}1\\-1\end{pmatrix}
  =-
  \begin{pmatrix}1\\-1\end{pmatrix}.
  \label{eq:l6-P45-check}
\end{equation}

\begin{classroomqa}
\audiencequestion{Does knowing the state determine what is
observable?\lecturetimestamp{00:35:14--00:37:32}}
\lecturerresponse{No. The state determines which observables, if any, have
definite values. In \(\ket{d_+}\), the diagonal observable is definite while
\(\hat P_0\) gives two outcomes with equal probability. Similarly, a
position delta state makes position definite, whereas a plane wave makes
momentum definite; a generic wavefunction usually makes neither definite.}
\end{classroomqa}

\section{Arbitrary Linear Polarization}

Let the transmission axis make angle \(\theta\) with the \(x\)-axis. Its
state and orthogonal partner are
\begin{align}
  \ket{\theta}
    &=\cos\theta\,\ket{x}+\sin\theta\,\ket{y}
      =
      \begin{pmatrix}
        \cos\theta\\
        \sin\theta
      \end{pmatrix},
      \label{eq:l6-theta-state}\\
  \ket{\theta_\perp}
    &=\ket{\theta+\pi/2}
      =-\sin\theta\,\ket{x}+\cos\theta\,\ket{y}
      =
      \begin{pmatrix}
        -\sin\theta\\
        \cos\theta
      \end{pmatrix}.
      \label{eq:l6-theta-perp}
\end{align}
Both have unit norm, and
\begin{equation}
  \braket{\theta}{\theta_\perp}
  =-\cos\theta\sin\theta+\sin\theta\cos\theta=0.
  \label{eq:l6-theta-orthogonality}
\end{equation}
Also \(\ket{\theta+\pi}=-\ket{\theta}\), confirming that a linear
polarization axis has period \(\pi\), not \(2\pi\).

\begin{figure}[tbp]
  \centering
  \includegraphics[width=0.62\textwidth]{lecture_06_figure_04.png}
  \caption{The arbitrary linear-polarization state and its orthogonal
  companion written in the axial basis.}
  \lectureframe{00:43:30}
\end{figure}

Suppose the first polarizer prepares \(\ket{\theta}\) and the second tests
the \(x\)-direction. Then
\begin{equation}
  \operatorname{Prob}(x|\theta)
  =|\braket{x}{\theta}|^2
  =\cos^2\theta,
  \qquad
  \operatorname{Prob}(y|\theta)=\sin^2\theta.
  \label{eq:l6-Malus-xy}
\end{equation}
Their sum is one. Reversing the inner product does not change the
probability because
\(\braket{\theta}{x}=\braket{x}{\theta}^{*}\).

For two arbitrary axes \(\alpha\) and \(\beta\), the needed trigonometric
identity is
\begin{equation}
  \cos(\alpha-\beta)
  =\cos\alpha\cos\beta+\sin\alpha\sin\beta.
  \label{eq:l6-cos-difference}
\end{equation}
The related sum, sine, and double-angle identities are useful enough to keep
beside it.

\begin{figure}[tbp]
  \centering
  \includegraphics[width=0.72\textwidth]{lecture_06_figure_05.png}
  \caption{The trigonometric identities used to compare two rotated
  polarization bases.}
  \lectureframe{00:54:30}
\end{figure}

The transition amplitude is
\begin{align}
  \braket{\beta}{\alpha}
  &=
  \begin{pmatrix}
    \cos\beta&\sin\beta
  \end{pmatrix}
  \begin{pmatrix}
    \cos\alpha\\
    \sin\alpha
  \end{pmatrix}\\
  &=\cos\alpha\cos\beta+\sin\alpha\sin\beta\\
  &=\cos(\alpha-\beta).
  \label{eq:l6-angle-amplitude}
\end{align}
Therefore the single-photon form of Malus's law is
\begin{equation}
  \boxed{\operatorname{Prob}(\beta|\alpha)
  =|\braket{\beta}{\alpha}|^2
  =\cos^2(\alpha-\beta).}
  \label{eq:l6-Malus}
\end{equation}

\begin{figure}[tbp]
  \centering
  \includegraphics[width=0.70\textwidth]{lecture_06_figure_06.png}
  \caption{The overlap is an amplitude; its absolute square is the
  transmission probability.}
  \lectureframe{00:57:30}
\end{figure}

Only the relative angle appears, as rotational symmetry requires. At equal
angles the probability is one; at orthogonal angles it is zero. For every
intermediate angle each photon still gives one of two discrete records.

\section{How an Extra Polarizer Opens a Closed Path}

Prepare \(\ket{\theta}\) and test it directly with the orthogonal analyzer
\(\theta+\pi/2\). The transmission probability vanishes:
\begin{equation}
  \operatorname{Prob}(\theta+\pi/2|\theta)
  =\cos^2(\pi/2)=0.
  \label{eq:l6-direct-block}
\end{equation}
Now insert an \(x\)-polarizer between them. Passing the intermediate
polarizer has probability \(\cos^2\theta\), and passage prepares the new
state \(\ket{x}\). Conditional on that event, transmission through the last
analyzer has probability
\begin{equation}
  \operatorname{Prob}(\theta+\pi/2|x)
  =\cos^2(\theta+\pi/2)=\sin^2\theta.
\end{equation}
The probability of completing both stages is therefore
\begin{equation}
  \operatorname{Prob}
  (\theta\longrightarrow x\longrightarrow\theta+\pi/2)
  =\cos^2\theta\,\sin^2\theta.
  \label{eq:l6-intermediate-polarizer}
\end{equation}

\begin{figure}[tbp]
  \centering
  \includegraphics[width=0.66\textwidth]{lecture_06_figure_07.png}
  \caption{An intermediate horizontal polarizer changes the state and
  produces the nonzero product
  \(\cos^2\theta\,\sin^2\theta\).}
  \lectureframe{01:02:30}
\end{figure}

\begin{classroomqa}
\audiencequestion{How can inserting another obstacle increase the chance of
transmission?\lecturetimestamp{00:59:29--01:03:40}}
\lecturerresponse{Without the middle apparatus the initial and final states
are orthogonal. The intermediate polarizer is not a passive filter on an
unchanged state: successful passage prepares \(\ket{x}\), which has a
nonzero overlap with the final state. The two conditional probabilities
multiply. Adding a measurement can therefore open a route that was absent
in the direct experiment.}
\end{classroomqa}

\begin{classroomqa}
\audiencequestion{Why can crossed polarizers show colors when stressed
transparent material is placed between them?
\lecturetimestamp{01:03:42--01:05:00}}
\lecturerresponse{The exchange in the lecture leaves the mechanism open and
notes that real optical elements can have wavelength-dependent efficiency.
More specifically, stressed transparent materials are often birefringent:
they give two polarization components a wavelength-dependent relative phase.
The final analyzer converts that phase delay into wavelength-dependent
intensity. The photon frequency is not changed by an ideal static
polarizer; different frequencies are transmitted with different
probabilities.}
\end{classroomqa}

\section{The Observable at Angle \texorpdfstring{\(\theta\)}{theta}}

The general analyzer must have \(\ket{\theta}\) and
\(\ket{\theta_\perp}\) as eigenvectors:
\begin{equation}
  \hat P_\theta\ket{\theta}=+\ket{\theta},
  \qquad
  \hat P_\theta\ket{\theta_\perp}=-\ket{\theta_\perp}.
  \label{eq:l6-Ptheta-eigenvalues}
\end{equation}
Its spectral form is already the answer:
\begin{equation}
  \hat P_\theta
  =\ket{\theta}\bra{\theta}
   -\ket{\theta_\perp}\bra{\theta_\perp}.
  \label{eq:l6-Ptheta-spectral}
\end{equation}
In the \(x,y\) basis,
\begin{align}
  \hat P_\theta
  &=
  \begin{pmatrix}
    \cos^2\theta-\sin^2\theta&
    2\sin\theta\cos\theta\\
    2\sin\theta\cos\theta&
    \sin^2\theta-\cos^2\theta
  \end{pmatrix}\\
  &=
  \begin{pmatrix}
    \cos2\theta&\sin2\theta\\
    \sin2\theta&-\cos2\theta
  \end{pmatrix}.
  \label{eq:l6-Ptheta-matrix}
\end{align}

Let \(c=\cos\theta\) and \(s=\sin\theta\). Acting on the proposed
\(+1\) eigenvector gives
\begin{align}
  \hat P_\theta
  \begin{pmatrix}c\\s\end{pmatrix}
  &=
  \begin{pmatrix}
    c^2-s^2&2sc\\
    2sc&s^2-c^2
  \end{pmatrix}
  \begin{pmatrix}c\\s\end{pmatrix}\\
  &=
  \begin{pmatrix}
    c(c^2+s^2)\\
    s(c^2+s^2)
  \end{pmatrix}
  =
  \begin{pmatrix}c\\s\end{pmatrix}.
  \label{eq:l6-Ptheta-check}
\end{align}

\begin{figure}[tbp]
  \centering
  \includegraphics[width=0.62\textwidth]{lecture_06_figure_08.png}
  \caption{The board check of the general polarization matrix on
  \((\cos\theta,\sin\theta)^T\).}
  \lectureframe{01:10:30}
\end{figure}

Rotating the analyzer by \(90^\circ\) exchanges the two outcomes:
\begin{equation}
  \hat P_{\theta+\pi/2}=-\hat P_\theta.
  \label{eq:l6-orthogonal-sign}
\end{equation}
The trigonometric reason is that both
\(\cos(2\theta+\pi)\) and \(\sin(2\theta+\pi)\) acquire a minus sign. The
physical reason is simpler: the state formerly labelled \(+1\) is now
labelled \(-1\), and conversely.

\begin{figure}[tbp]
  \centering
  \includegraphics[width=0.58\textwidth]{lecture_06_figure_09.png}
  \caption{The two orthogonal analyzer axes whose outcome labels are
  exchanged by a \(90^\circ\) rotation.}
  \lectureframe{01:15:15}
\end{figure}

\section{Complex Amplitudes and Circular Polarization}

Real coefficients describe every linear polarization, but not every vector
in a complex two-dimensional Hilbert space. Consider
\begin{equation}
  \ket{c_+}
  =\frac{1}{\sqrt2}
    \begin{pmatrix}1\\i\end{pmatrix},
  \qquad
  \ket{c_-}
  =\frac{1}{\sqrt2}
    \begin{pmatrix}1\\-i\end{pmatrix}.
  \label{eq:l6-circular-states}
\end{equation}
Complex conjugation is essential:
\begin{equation}
  \braket{c_+}{c_+}
  =\frac12\bigl(1+(-i)i\bigr)=1,
\end{equation}
and
\begin{equation}
  \braket{c_-}{c_+}
  =\frac12\bigl(1+i^2\bigr)=0.
  \label{eq:l6-circular-orthogonality}
\end{equation}
The tempting product without conjugating the bra gives the wrong answer.

\begin{figure}[tbp]
  \centering
  \includegraphics[width=0.60\textwidth]{lecture_06_figure_10.png}
  \caption{The two normalized complex vectors that form the circular
  polarization basis.}
  \lectureframe{01:21:00}
\end{figure}

To see their wave meaning, restore a dimensionally explicit phase
\(\phi=kz-\omega t\), with \(\omega=ck\). Equal in-phase components,
\begin{equation}
  E_x=E_0\cos\phi,
  \qquad
  E_y=E_0\cos\phi,
\end{equation}
make a linearly polarized wave at \(45^\circ\). A relative quarter-cycle
shift gives
\begin{equation}
  E_x=E_0\sin\phi,
  \qquad
  E_y=E_0\cos\phi.
  \label{eq:l6-circular-field}
\end{equation}
At a fixed point in space, the tip of
\(\bm E=E_x\hat{\bm x}+E_y\hat{\bm y}\) moves around a circle. Reversing the
sign of one quadrature reverses the sense of rotation. The magnetic field
is transverse as well; in vacuum,
\(\bm B=\hat{\bm z}\times\bm E/c\).

\begin{figure}[tbp]
  \centering
  \includegraphics[width=0.68\textwidth]{lecture_06_figure_11.png}
  \caption{The \(x\)- and \(y\)-components are displaced by one quarter
  cycle.}
  \lectureframe{01:24:30}
\end{figure}

\begin{figure}[tbp]
  \centering
  \includegraphics[width=0.54\textwidth]{lecture_06_figure_12.png}
  \caption{At fixed \(z\), the two quadratures make the electric-field
  vector rotate.}
  \lectureframe{01:25:30}
\end{figure}

Which state is called right- or left-circular depends on the propagation,
time, and viewing conventions. The invariant content here is the orthogonal
pair \((1,\pm i)^T/\sqrt2\).

\begin{classroomqa}
\audiencequestion{What optical device creates the relative phase required
for circular polarization?\lecturetimestamp{01:26:37--01:27:05}}
\lecturerresponse{A quarter-wave plate. Its two principal axes accumulate
different phases; a linear state with equal components along those axes
emerges with a relative phase of \(\pi/2\), producing circular
polarization. Unequal component magnitudes generally produce an ellipse.}
\end{classroomqa}

For any linear analyzer,
\begin{equation}
  \braket{\theta}{c_+}
  =\frac{1}{\sqrt2}
    \bigl(\cos\theta+i\sin\theta\bigr).
\end{equation}
Hence
\begin{align}
  |\braket{\theta}{c_+}|^2
  &=
  \frac12
  (\cos\theta+i\sin\theta)
  (\cos\theta-i\sin\theta)\\
  &=\frac12(\cos^2\theta+\sin^2\theta)
  =\frac12.
  \label{eq:l6-circular-linear-probability}
\end{align}
A circularly polarized photon therefore passes every linear analyzer with
probability one half. Its statistics reveal no preferred linear axis.

\begin{figure}[tbp]
  \centering
  \includegraphics[width=0.70\textwidth]{lecture_06_figure_13.png}
  \caption{The complex-conjugate product removes the phase and leaves
  transmission probability \(1/2\).}
  \lectureframe{01:29:50}
\end{figure}

The general normalized pair \(a,b\in\mathbb C\), modulo a common phase,
describes elliptical polarization. A real relative phase gives linear
polarization. Equal magnitudes with relative phase \(\pm\pi/2\) give the two
circular polarizations. Intermediate magnitudes and phases give ellipses.

\begin{classroomqa}
\audiencequestion{What observable distinguishes the two circular states?
\lecturetimestamp{01:32:19--01:32:36}}
\lecturerresponse{The lecture leaves this construction as an exercise. With
\(\ket{c_+}\) assigned eigenvalue \(+1\) and \(\ket{c_-}\) assigned
\(-1\), the spectral rule gives
\[
  \hat P_c
  =\ket{c_+}\bra{c_+}-\ket{c_-}\bra{c_-}
  =
  \begin{pmatrix}
    0&-i\\
    i&0
  \end{pmatrix}.
\]
Changing the handedness convention or exchanging the labels multiplies the
operator by \(-1\).}
\end{classroomqa}

\section{Compatible and Incompatible Polarization Tests}

The axial, diagonal, arbitrary linear, and circular observables do not
usually share eigenvectors. In the present two-state space,
\begin{equation}
  \hat P_\theta
  =\cos2\theta\,\sigma_z+\sin2\theta\,\sigma_x,
\end{equation}
so two linear analyzers satisfy
\begin{equation}
  [\hat P_\alpha,\hat P_\beta]
  =2i\sin\!\bigl(2(\beta-\alpha)\bigr)\sigma_y.
  \label{eq:l6-linear-commutator}
\end{equation}
They commute only when their axes are the same or orthogonal, modulo the
\(\pi\)-periodicity of a linear polarization axis. In the orthogonal case
\(\hat P_{\theta+\pi/2}=-\hat P_\theta\), so measuring one immediately fixes
the other with the opposite label.

\begin{classroomqa}
\audiencequestion{Can \(x\)- and \(y\)-polarization be measured
simultaneously, and what changes at \(45^\circ\)?
\lecturetimestamp{01:33:31--01:35:20}}
\lecturerresponse{The \(x\) and \(y\) outcome labels belong to one common
basis and are negatives of one another in the \(\pm1\) convention. Knowing
one fixes the other. The \(x\)- and \(45^\circ\)-analyzers are genuinely
different: their operators do not commute and have no common eigenstate.
The same incompatibility holds between a linear analyzer and the circular
observable.}
\end{classroomqa}

Incompatibility is therefore not vagueness about an apparatus. Each
apparatus is perfectly definite. It is the impossibility of one state being
a simultaneous eigenstate of two noncommuting tests.

\section{Expectation Value Means Average Value}

The standard term \emph{expectation value} can be misleading. Imagine a
position probability density concentrated in two symmetric, separated
lobes. Its average position is zero, even if the probability of finding the
particle near zero is tiny or vanishes. The average need not be a likely
single outcome.

\begin{figure}[tbp]
  \centering
  \includegraphics[width=0.62\textwidth]{lecture_06_figure_14.png}
  \caption{A symmetric two-lobe profile can have average position zero even
  though its central probability is very small.}
  \lectureframe{01:36:09}
\end{figure}

Let \(\hat K\) be a Hermitian observable with a complete orthonormal
eigenbasis,
\begin{equation}
  \hat K\ket{n}=\lambda_n\ket{n},
  \qquad
  \sum_n\ket{n}\bra{n}=\mathbf 1.
  \label{eq:l6-K-spectrum}
\end{equation}
For a normalized state \(\ket{\psi}\), define
\begin{equation}
  \expect{\hat K}_\psi
  :=\bra{\psi}\hat K\ket{\psi}.
  \label{eq:l6-expectation-definition}
\end{equation}

\begin{classroomqa}
\audiencequestion{Should there be a complex-conjugation star on
\(\hat K\)?\lecturetimestamp{01:38:17--01:38:41}}
\lecturerresponse{No additional star is needed. Turning
\(\ket{\psi}\) into \(\bra{\psi}\) already takes the complex conjugate
transpose. The observable itself is Hermitian,
\(\hat K^\dagger=\hat K\), which ensures that the expectation value is
real.}
\end{classroomqa}

Insert the eigenbasis resolution of identity between \(\hat K\) and
\(\ket{\psi}\):
\begin{align}
  \expect{\hat K}_\psi
  &=
  \bra{\psi}\hat K
  \left(\sum_n\ket{n}\bra{n}\right)
  \ket{\psi}\\
  &=
  \sum_n
  \bra{\psi}\hat K\ket{n}
  \braket{n}{\psi}\\
  &=
  \sum_n
  \lambda_n
  \braket{\psi}{n}
  \braket{n}{\psi}\\
  &=
  \sum_n\lambda_n
  |\braket{n}{\psi}|^2.
  \label{eq:l6-expectation-spectral}
\end{align}

\begin{figure}[tbp]
  \centering
  \includegraphics[width=0.62\textwidth]{lecture_06_figure_15.png}
  \caption{The expectation-value sandwich after inserting the eigenbasis
  resolution of identity.}
  \lectureframe{01:39:45}
\end{figure}

The factor
\begin{equation}
  p_n=|\braket{n}{\psi}|^2
\end{equation}
is the probability of obtaining \(\lambda_n\). Thus
\(\expect{\hat K}_\psi=\sum_n p_n\lambda_n\) is exactly an ordinary weighted
average of the possible outcomes.

\begin{figure}[tbp]
  \centering
  \includegraphics[width=0.60\textwidth]{lecture_06_figure_16.png}
  \caption{Possible eigenvalues weighted by their probabilities give the
  average \(\sum_n p_n\lambda_n\).}
  \lectureframe{01:42:50}
\end{figure}

Return finally to the axial polarization observable in the state
\(\ket{\theta}\):
\begin{align}
  \expect{\hat P_0}_\theta
  &=
  \bra{\theta}
  \begin{pmatrix}
    1&0\\
    0&-1
  \end{pmatrix}
  \ket{\theta}\\
  &=
  \begin{pmatrix}
    \cos\theta&\sin\theta
  \end{pmatrix}
  \begin{pmatrix}
    \cos\theta\\
    -\sin\theta
  \end{pmatrix}\\
  &=
  \cos^2\theta-\sin^2\theta
  =\cos2\theta.
  \label{eq:l6-P0-expectation}
\end{align}

\begin{figure}[H]
  \centering
  \includegraphics[width=0.52\textwidth]{lecture_06_figure_17.png}
  \caption{The axial polarization expectation value in the state
  \(\ket{\theta}\).}
  \lectureframe{01:45:50}
\end{figure}

The same answer follows directly from the two outcomes:
\begin{equation}
  (+1)\cos^2\theta+(-1)\sin^2\theta=\cos2\theta.
\end{equation}
No single detector click reports \(\cos2\theta\). Each run still records
\(+1\) or \(-1\); \(\cos2\theta\) emerges only as the average over repeated
identically prepared photons. That distinction between a state, a
single-shot outcome, a probability, and an ensemble average completes the
two-state polarization example.
